\section{Related Work}
\label{sec:related}

\paragraph{Linear probing in language models.}
Linear probing of intermediate representations is a well-established
interpretability tool \cite{alain2016understanding, hewitt2019structural,
belinkov2022probing}. Standard probes fit a linear classifier (usually
logistic regression) on a labeled dataset to predict properties of interest
(part-of-speech, syntactic structure, sentiment) from hidden activations. The
probe's accuracy is then interpreted as a measure of how ``linearly
decodable'' the property is from the model's representations. Our method
differs in three ways: (i)~we use \emph{zero} labeled data---only LLM-generated
stories with pseudo-labels inherited from the generation prompt;
(ii)~we use \emph{unsupervised} PCA of nine centroids rather than supervised
logistic regression; (iii)~we are specifically interested in the
\emph{direction} of variation, not just the decodability.

\paragraph{Concept direction discovery and representation engineering.}
A parallel line of work extracts concept directions for activation steering
and interpretability
\cite{turner2024activation, rimsky2024steering, zou2023representation,
arditi2024refusal}. These methods typically use contrastive pairs: pairs
of prompts differing in the target concept (honest/dishonest,
harmful/harmless, happy/sad). The direction is the mean-difference of the
paired activations at a chosen layer. Park, Choe and Veitch
\cite{park2024linear} formalise the assumption that such linear directions
exist as the \emph{linear representation hypothesis} and derive geometric
tests (principal angle, subspace alignment) for when a concept lies in a
low-dimensional affine slice. \ours{} differs from this family in two
practical ways. First, the operator does \emph{not} write contrastive
pairs; we sample nine stories per category and let the first principal
component of the nine centroids discover the polarity axis automatically.
Second, we compare directly: on matched Qwen-1.5B L28 features, the
mean-difference contrastive axis gives SST-2 AUC $0.753$, while
PCA-of-nine-centroids gives $0.868$ ($+0.115$). The 9-centroid geometry
therefore recovers a stronger polarity direction than the binarised
contrastive alternative, while requiring less operator effort.
Relatedly, linear-direction work on lexical polarity
\cite{mikolov2013linguistic, bolukbasi2016man} establishes that
linguistically-meaningful concepts can live on single directions in
pretrained embeddings; \ours{} extends this to \emph{instruction-tuned
LLM residual streams} and to \emph{sentence-level} valence with zero
labels.

\paragraph{Emotion-vector distillation from LLMs.}
Several recent papers use LLM-derived emotion vectors as auxiliary signals
for downstream tasks \cite{zhang2024emotionkd, li2025emod, aqa2024emotion}.
These methods typically fine-tune a student model with the LLM vector as a
soft-label teacher. Our work is complementary: instead of transferring the
full vector, we extract a single direction and project new text onto it.
We also show that for classification (our downstream tasks), the LLM's
content is \emph{not the source of the signal}---a random 9-dimensional
orthonormal prototype tensor matches the LLM-derived one within noise.

\paragraph{Sentence embeddings and sentiment lexicons.}
Sentence encoders such as SBERT \cite{reimers2019sentence} and SimCSE
\cite{gao2021simcse} provide fixed-length text representations that can be
used for similarity search. We compare directly to SBERT
(\texttt{all-MiniLM-L6-v2}) with the same nine-story centroid pipeline and
find that the generative LLM's residual stream is \emph{substantially} richer
for this purpose (+15 to +24 percentage points across SST-2, IMDB, and Hu \&
Liu). Classical sentiment lexicons such as NRC VAD
\cite{mohammad2018obtaining}, Warriner \textit{et al.\@}
\cite{warriner2013norms}, and Hu \& Liu \cite{huliu2004opinion} require
thousands to tens of thousands of hand-labeled words. We reach competitive
correlation with these lexicons (Pearson 0.62--0.76) using only
the nine story centroids and no human labels.

\paragraph{Cross-lingual representation transfer.}
Multilingual sentence encoders \cite{conneau2020unsupervised} share a
representation space across languages by explicit multilingual
pre-training. Instruction-tuned LLMs are trained on large mostly-English
corpora with smaller non-English sub-corpora, and it is not guaranteed
that concept directions extracted from English internals will transfer to
other languages. We show that the English V-axis, applied zero-shot to
held-out stories in Spanish, French, German, and Chinese, produces AUCs
in the $0.949$--$1.000$ range, and that within-language extracted axes
are similarly usable. This is a practical cross-lingual robustness result
complementing multilingual encoder work: a single zero-label English
axis is a usable multilingual valence probe on the multilingual
instruction-tuned LLM we tested.

\paragraph{Sudden-emergence findings.}
Phase-transition phenomena in transformer training have received recent
attention \cite{olsson2022induction, wei2022emergent, michaud2024quantization}.
Most of these papers study training dynamics (abilities emerging at particular
scale or steps). We report an emergence \emph{in layer space} at inference: the
valence direction appears abruptly at layer 27/28, not gradually through the
network. To our knowledge this abrupt-in-depth pattern has not been reported;
standard probing papers test one or two layers and miss the V-shape.

\paragraph{Few-shot prompting baselines.}
Direct-prompting sentiment classification with instruction-tuned LLMs is
strong \cite{brown2020language, wei2022finetuned}. We benchmark against this
baseline and find that prompting beats probing on accuracy (92.78\% vs.\
77.75\% on SST-2). Our method's advantages are therefore \emph{not}
classification accuracy but (i)~\speedup{} inference speedup via batched
forward passes, (ii)~continuous scores useful for ranking, calibration, and
distribution monitoring, (iii)~interpretability insights about where and how
the LLM represents the concept, and (iv)~no prompt engineering.
